\subsection{Gubaev et al.\ (2023): MTP for the TaVCrW high-entropy alloy}
\label{sec:appendix-example-gubaev2023}

\paragraph{Q1 -- Bibliographic identification.}
Gubaev, Zaverkin, Srinivasan, Duff, K\"astner, and Grabowski compare
two complementary machine-learned potentials---a Moment Tensor
Potential (MTP) and a Gaussian-Moment Neural Network (GM-NN)---on the
TaVCrW BCC refractory high-entropy alloy. Published in \emph{npj
Computational Materials} 9, 129 (2023);
doi:10.1038/s41524-023-01073-w.
\lstinputlisting[language=turtle]{artifacts/kg/listings/gubaev2023/q01-bibliographic.ttl}

\paragraph{Q2 -- Material system.}
The studied system is the equiatomic TaVCrW BCC refractory high-entropy
alloy (space group $Im\bar{3}m$). Training and validation sets span
its 2-, 3-, and 4-component subsystems, plus a deformed B2/B2
phase-segregated quaternary used as out-of-distribution test. We
model the chemistry at the four-component level and record the
subsystem coverage in $\dlRole{materialClass}$.
\lstinputlisting[language=turtle]{artifacts/kg/listings/gubaev2023/q02-material.ttl}

\paragraph{Q3 -- Reference calculation method.}
Reference data are produced from DFT calculations.
\lstinputlisting[language=turtle]{artifacts/kg/listings/gubaev2023/q03-reference-method-type.ttl}

\paragraph{Q4 -- Reference settings.}
DFT calculations are performed with VASP at the PBE/GGA level using
PAW pseudopotentials (11 valence electrons for Ta and V; 12 for Cr
and W; spin-unpolarised), with Methfessel--Paxton smearing of order 1
($\sigma=0.1$), a 350\,eV plane-wave cutoff, and an automatic k-mesh
generated by VASP at $\text{KSPACING}=0.12$.
\lstinputlisting[language=turtle]{artifacts/kg/listings/gubaev2023/q04-reference-settings.ttl}

\paragraph{Q5 -- Training dataset.}
The training+validation set has 6{,}711 DFT configurations: 5{,}680 at
0\,K (4{,}491 binary, 595 ternary, 594 quaternary) plus 1{,}031
MD-sampled at 2{,}500\,K (close to the melting point). Energies,
forces, and stresses are covered. Provenance is \emph{in-house}.
\lstinputlisting[language=turtle]{artifacts/kg/listings/gubaev2023/q05-dataset.ttl}

\paragraph{Q6 -- Sampling strategies.}
Four strategies are combined: 0\,K small-supercell enumeration with
\texttt{enumlib} (2--8 atom non-repetitive symmetry-unique supercells;
2{,}500 binary structures plus ternaries, equiatomic and off-equiatomic
quaternaries); vibrational sampling via DFT MD at 2{,}500\,K; active
learning using the MaxVol extrapolation grade (selecting MD frames
with grade above threshold for DFT recomputation); and
phase-segregation sampling via 432-atom B2/B2 deformed-binary
structures.
\lstinputlisting[language=turtle]{artifacts/kg/listings/gubaev2023/q06-sampling.ttl}

\paragraph{Q7 -- MLIP method, components, implementation.}
We encode the MTP arm. The method is a Moment Tensor Potential with a
level-restricted basis ($\mathit{lev}_{\max}$ controls the radial+angular
truncation), implemented in the MLIP package
(\url{https://gitlab.com/ashapeev/mlip-2}). Loss is weighted MSE on
energies/forces/stresses with weights
$w_E=1/N_{\text{at}}$, $w_f=0.1$, $w_\sigma=0.001$.
\lstinputlisting[language=turtle]{artifacts/kg/listings/gubaev2023/q07-method.ttl}

\paragraph{Q8 -- Hyperparameter settings.}
Reported setting: cutoff radius $r_c=5.0$\,\AA. The
$\mathit{lev}_{\max}$ value is referenced under
$\dlRole{hasHyperparameter}$ but the paper compares several values
(0, 2, $\dots$) without designating one as canonical, so we omit a
$\HyperparameterSettingC$ instance for it.
\lstinputlisting[language=turtle]{artifacts/kg/listings/gubaev2023/q08-hyperparameter-settings.ttl}

\paragraph{Q9 -- MLIP run and trained model.}
A single MTP training run on the 6{,}711-configuration training set
produces one trained TaVCrW MTP. No active-learning ensemble
structure is recorded at the level of separate
$\dlConcept{MLIPRun}$ instances.
\lstinputlisting[language=turtle]{artifacts/kg/listings/gubaev2023/q09-run-and-model.ttl}

\paragraph{Q10 -- Benchmark results.}
Overall RMSEs across the in-distribution Ta--V--Cr--W subsystems
(Table 1, MTP): 2.43\,meV/atom on energy, 0.054\,eV/\AA{} on
forces. The 2{,}500\,K disordered TaVCrW point alone has RMSEs of
2.40\,meV/atom and 0.156\,eV/\AA{}; the cross-method comparison with
GM-NN and EAM is in the same table.
\lstinputlisting[language=turtle]{artifacts/kg/listings/gubaev2023/q10-benchmark-results.ttl}

\paragraph{Q11 -- Computational resources.}
\emph{Not reported} in absolute form. The paper compares MTP and
GM-NN convergence with training-set size and execution speed
qualitatively but gives no training duration, GPU/CPU hours,
training hardware, peak memory, or per-atom inference time.
\lstinputlisting[language=turtle]{artifacts/kg/listings/gubaev2023/q11-resources.ttl}

\paragraph{Q12 -- Gaps.}
The most consequential omission is the GM-NN arm: the paper trains
\emph{both} an MTP and a GM-NN on the same data and frames the
comparison as the main contribution. The ontology can express each
as an $\dlConcept{MLIPMethod}$ instance with its own
$\dlConcept{MLIPRun}$ and $\dlConcept{TrainedModel}$, and a sibling
\texttt{gubaev2023-gmnn.ttl} would mirror this file with
$\texttt{ex:GM-NN}$ and the \texttt{gm-nn} package as
$\dlConcept{Library}$. Other unreported items: $\mathit{lev}_{\max}$
as a $\HyperparameterSettingC$ (the paper sweeps several values);
the GM-NN architectural hyperparameters (number of layers, channel
widths, learning rate); active-learning iteration counts; downstream
phonon and elastic-constant benchmarks; and compute-cost metadata.
